\uuid{UHqZ}
\titre{Pour tout/il existe}
\auteur{bodin ; hanani ; mzari}
\organisation{exo7}
\licence{CC BY-NC-SA 4.0}
\datecreate{2020-01-20}
\niveau{L1}
\module{Algèbre}
\chapitre{Logique, ensemble, raisonnement}
\sousChapitre{Logique}
\difficulte{1}
\theme{Logique, pour tout, il existe}
\exo7qcmid{lille-1.logique.005}

\contenu{
\begin{question}
\qclassification{100.01}
\qtags{motcle=pour tout/il existe}
Un groupe de coureurs $C$ chronomètre ses temps : $t(c)$ désigne le temps (en secondes) du coureur $c$.
Dans ce groupe Valentin et Chloé ont réalisé le meilleur temps de $47$ secondes. Tom est déçu car il est arrivé troisième, avec un temps de $55$ secondes.
À partir de ces informations, quelles sont les assertions dont on peut déduire qu'elles sont vraies ?
\begin{answers}
    \good{$\forall c \in C \quad t(c) \ge 47$}

    \bad{$\exists c \in C \quad 47 < t(c) < 55$}

    \good{$\exists c \in C \quad t(c) > 47$}

    \bad{$\forall c \in C \quad t(c) \le 55$}
\end{answers}
\begin{explanations}
Comme Tom est troisième, il n'existe pas de $c$ tel que $47 < t(c) < 55$.
\end{explanations}
\end{question}
}
